\section*{Chapter \ref{ch:g11:functional-groups} --- Functional Groups and Families}

\begin{solution}{exo:g11:functional-groups:1}
Hydroxyl, \omterm{def:g11:functional-groups:alcohol}{alcohol}; carbonyl inside the chain, \omterm{def:g11:functional-groups:carbonyl}{ketone}; carboxyl,
\omterm{def:g11:functional-groups:carboxylic-acid}{carboxylic acid}; \omterm{def:g11:functional-groups:carboxylic-acid}{ester} group, \omterm{def:g11:functional-groups:carboxylic-acid}{ester}; amino group \ce{-NH2}, \omterm{def:g11:functional-groups:amine}{amine}.
\end{solution}

\begin{solution}{exo:g11:functional-groups:2}
Methanol; propan-1-ol; methanoic acid; propanal; propanone.
\end{solution}

\begin{solution}{exo:g11:functional-groups:3}
\omterm{def:g11:functional-groups:alcohol}{Alcohol}; \omterm{def:g11:functional-groups:carbonyl}{aldehyde}; \omterm{def:g11:functional-groups:carboxylic-acid}{ester}; \omterm{def:g11:functional-groups:amine}{amide}; \omterm{def:g11:functional-groups:halogenoalkane}{halogenoalkane}; \omterm{def:g11:functional-groups:halogenoalkane}{alkene}.
\end{solution}

\begin{solution}{exo:g11:functional-groups:4}
Propanal \ce{CH3-CH2-CHO}, propanone \ce{CH3-CO-CH3}. Propanal is the
\omterm{def:g11:functional-groups:carbonyl}{aldehyde}: its \ce{C=O} carbon is at the end of the chain, bonded to a
hydrogen \omterm{def:g7:atoms-and-molecules:atom}{atom}; in propanone it is between two carbon \omterm{def:g7:atoms-and-molecules:atom}{atoms}.
\end{solution}

\begin{solution}{exo:g11:functional-groups:5}
Primary (the carbon bearing \ce{-OH} has one carbon neighbour);
secondary (two); tertiary (three).
\end{solution}

\begin{solution}{exo:g11:functional-groups:6}
\setchemfig{atom sep=1.5em}
Pentan-3-ol \chemfig{-[:30]-[:-30](-[:-90]OH)-[:30]-[:-30]};
2-methylbutanoic acid \chemfig{-[:30]-[:-30](-[:-90])-[:30](=[:90]O)-[:-30]OH};
hexan-2-one \chemfig{-[:30](=[:90]O)-[:-30]-[:30]-[:-30]-[:30]};
propyl ethanoate \chemfig{-[:30](=[:90]O)-[:-30]O-[:30]-[:-30]-[:30]}.
\end{solution}

\begin{solution}{exo:g11:functional-groups:7}
\ce{C3H6O}. Yes, they are \omterm{def:g11:organic-skeletons:isomer}{isomers}, but with different groups: an
\omterm{def:g11:functional-groups:carbonyl}{aldehyde} and a \omterm{def:g11:functional-groups:carbonyl}{ketone} (both carbonyl groups, in different places).
\end{solution}

\begin{solution}{exo:g11:functional-groups:8}
\ce{HCOO-CH2-CH3}, ethyl methanoate; \ce{CH3-COO-CH2-CH2-CH3}, propyl
ethanoate.
\end{solution}

\begin{solution}{exo:g11:functional-groups:9}
Methyl ethanoate (ethanoic acid, methanol); ethyl propanoate (propanoic
acid, ethanol); methyl methanoate (methanoic acid, methanol).
\end{solution}

\begin{solution}{exo:g11:functional-groups:10}
\setchemfig{atom sep=1.5em}
\begin{center}
\small
\begin{tabular}{cccc}
\chemfig{-[:30]-[:-30]-[:30]-[:-30]OH} &
\chemfig{-[:30](-[:90]OH)-[:-30]-[:30]} &
\chemfig{-[:30](-[:90])-[:-30]-[:30]OH} &
\chemfig{-[:0](-[:90]OH)(-[:-90])-[:0]} \\
butan-1-ol & butan-2-ol & 2-methylpropan-1-ol & 2-methylpropan-2-ol \\
primary & secondary & primary & tertiary
\end{tabular}
\end{center}
\end{solution}

\begin{solution}{exo:g11:functional-groups:11}
An \omterm{def:g11:functional-groups:amine}{amide} (ethanamide). \emph{-ene}. \omterm{def:g11:functional-groups:carboxylic-acid}{Carboxylic acids} (\ce{C=O} and
\ce{-OH} on the same carbon), \omterm{def:g11:functional-groups:carboxylic-acid}{esters} (\ce{C=O} and \ce{-O-}) and \omterm{def:g11:functional-groups:amine}{amides}
(\ce{C=O} and \ce{N}): three families.
\end{solution}

\begin{solution}{exo:g11:functional-groups:12}
Aspirin: a carboxyl group (\omterm{def:g11:functional-groups:carboxylic-acid}{carboxylic acid}) and an \omterm{def:g11:functional-groups:carboxylic-acid}{ester} group.
Paracetamol: a phenol \ce{-OH} on the ring and an \omterm{def:g11:functional-groups:amine}{amide} group
\ce{NH-CO}.
\end{solution}

\begin{solution}{exo:g11:functional-groups:13}
\ce{C2H6O}. Only ethanol, with its \ce{O-H}, forms \omterm{def:g11:polarity-and-cohesion:hydrogen-bond}{hydrogen bonds}
between its \omterm{def:g7:atoms-and-molecules:molecule}{molecules} (methoxymethane has no hydrogen on its oxygen).
Ethanol boils higher, at \qty{78}{\celsius}; methoxymethane at
\qty{-25}{\celsius}.
\end{solution}

\begin{solution}{exo:g11:functional-groups:14}
A hydroxyl group and a carboxyl group. The chain has three carbons, the
acid carbon being carbon 1 and the \ce{-OH} on carbon 2:
2-hydroxypropanoic acid. Glycine: 2-aminoethanoic acid (often written
aminoethanoic acid).
\end{solution}

\begin{solution}{exo:g11:functional-groups:15}
Ethanol \ce{CH3-CH2-OH}, \ce{C2H6O}, \omterm{def:g11:functional-groups:alcohol}{alcohol}; ethanal \ce{CH3-CHO},
\ce{C2H4O}, \omterm{def:g11:functional-groups:carbonyl}{aldehyde}; ethanoic acid \ce{CH3-COOH}, \ce{C2H4O2},
\omterm{def:g11:functional-groups:carboxylic-acid}{carboxylic acid}. From ethanol to ethanal two hydrogen \omterm{def:g7:atoms-and-molecules:atom}{atoms} are lost;
from ethanal to ethanoic acid one oxygen \omterm{def:g7:atoms-and-molecules:atom}{atom} is gained.
\end{solution}

\begin{solution}{pb:g11:functional-groups:1}
\textbf{1.} A: \omterm{def:g11:functional-groups:carboxylic-acid}{ester} group; B: hydroxyl; C: carboxyl; D: carbonyl at the
end of the chain; E: \omterm{def:g11:functional-groups:carboxylic-acid}{ester} group.

\textbf{2.} A \omterm{def:g11:functional-groups:carboxylic-acid}{ester}, B \omterm{def:g11:functional-groups:alcohol}{alcohol}, C \omterm{def:g11:functional-groups:carboxylic-acid}{carboxylic acid}, D \omterm{def:g11:functional-groups:carbonyl}{aldehyde}, E \omterm{def:g11:functional-groups:carboxylic-acid}{ester}.

\textbf{3.} A (banana) and E (pineapple): fruity smells.

\textbf{4.} Primary: the carbon bearing \ce{-OH} is bonded to one carbon.

\textbf{5.} \ce{C6H12O}; for example hexan-2-one,
\ce{CH3-CO-CH2-CH2-CH2-CH3}.

\textbf{6.} The chain holding the \ce{C-OH} has four carbons, numbered
from the \ce{-OH}, with a methyl on carbon 3: 3-methylbutan-1-ol.

\textbf{7.} C: ethanoic acid; D: hexanal.

\textbf{8.} Ethyl butanoate, related to butanoic acid and ethanol.

\textbf{9.} \emph{Ethanoate}: the acid part \ce{CH3-CO}, from ethanoic
acid; \emph{3-methylbutyl}: the \omterm{def:g11:functional-groups:alcohol}{alcohol} part, from 3-methylbutan-1-ol.

\textbf{10.} C (ethanoic acid) and B (3-methylbutan-1-ol).

\textbf{11.} A \ce{C7H14O2}; B \ce{C5H12O}; C \ce{C2H4O2}.

\textbf{12.} $\ce{C2H4O2} + \ce{C5H12O}$ contains \ce{C7H16O3}, one
\ce{H2O} more than \ce{C7H14O2}: water.

\textbf{13.} \ce{C2H4O2 + C5H12O -> C7H14O2 + H2O}: C 7 = 7, H 16 = 16,
O 3 = 3.

\textbf{14.} $M(\text{B}) = 5 \times 12.0 + 12 \times 1.0 + 16.0 =
\qty{88.0}{g/mol}$; $M(\text{C}) = \qty{60.0}{g/mol}$.

\textbf{15.} $60.0 + 88.0 - 18.0 = \qty{130.0}{g/mol}$.

\textbf{16.} $7 \times 12.0 + 14 \times 1.0 + 2 \times 16.0 =
\qty{130.0}{g/mol}$: the same.

\textbf{17.} E, \ce{C6H12O2}: $6 \times 12.0 + 12 \times 1.0 + 2 \times
16.0 = \qty{116.0}{g/mol}$.

\textbf{18.} \ce{C7H14O2}: same formula as A, so an isomer, with the
same \omterm{def:g10:the-mole:molar-mass}{molar mass}. A \omterm{def:g10:the-mole:molar-mass}{molar mass} alone cannot tell \omterm{def:g11:organic-skeletons:isomer}{isomers} apart.

\textbf{19.} A and heptanoic acid have the same \omterm{def:g7:atoms-and-molecules:atom}{atoms} but different
groups, an \omterm{def:g11:functional-groups:carboxylic-acid}{ester} and a \omterm{def:g11:functional-groups:carboxylic-acid}{carboxylic acid}: one smells of banana, the other
rancid. The group, and where it sits, decides.

\textbf{20.} \qty{130.0}{g/mol}.
\end{solution}
